% Original question wording: NMES Assignment 2026, Sections 1 and 2.
% Line-wrap hyphenation is removed; mathematical notation is typeset in LaTeX.
\usepackage{tcolorbox}
\newcommand{\assignmentquestion}[1]{%
  \begin{tcolorbox}[colback=accent!4,colframe=accent!65,
    boxrule=0pt,leftrule=2pt,arc=0pt,boxsep=0pt,
    left=7pt,right=7pt,top=5pt,bottom=5pt,
    before skip=0pt,after skip=8pt,fontupper=\small]
  \textbf{Assignment question.} #1
  \end{tcolorbox}%
}
\newcommand{\questionfiles}[1]{%
  \par\smallskip{\footnotesize\color{accent}
  \textbf{Files / results (relative to Source-code).} #1\par}\medskip
}
\newcommand{\QuestionOneOne}{%
Describe your cable: construction, dimensions, materials, rated voltage, rated
current (ampacity) and the laying conditions for which that rating is given
(burial depth, soil thermal resistivity, ambient temperature, arrangement).
Give a table of electromagnetic and thermal properties with a source for each value.}
\newcommand{\QuestionOneTwo}{%
Choose the environment consistently with the datasheet rating conditions,
and justify its properties.}
\newcommand{\QuestionOneThree}{%
State the sheath bonding you assume: single-point bonded (no net sheath current)
or solidly bonded (circulating sheath currents allowed), and how the armour
and pipe are connected. This choice strongly affects the losses in Section 3.}
\newcommand{\QuestionOneFour}{%
State the excitation you impose: phase-to-ground voltage and phase current,
as peak phasors, with the phase sequence (see the conventions in Appendix A).}
\newcommand{\QuestionTwoOne}{%
Define your computational domain, the equations solved and the boundary
conditions. Where do you truncate the domain, and why? Discuss the role of
the grounded screens: which regions actually influence the field in the insulation?}
\newcommand{\QuestionTwoTwo}{%
Can you simplify the geometry of your cable without degrading the precision
(e.g. disregard details in the geometrical description)? Quantify the effect
of each simplification.}
\newcommand{\QuestionTwoThree}{%
How do the results vary with the mesh refinement? Plot the maximum field in
the insulation and the capacitance versus the number of degrees of freedom,
and identify where local refinement is needed.}
\newcommand{\QuestionTwoFour}{%
Define the admissible field strength at the conductor screen and at the core
screen from the datasheet (or, if absent, from a cited reference).
Is the healthy cable within these limits?}
\newcommand{\QuestionTwoFive}{%
Add a defect to the insulation of your cable. State its nature (e.g. gas void,
water tree, protrusion of the semiconducting screen), its properties, size and
position, and discuss what a defect in a 2D cross-section represents in 3D.
Compare the field with and without defect along a radial cut.
Does the cable withstand the defect?}
\newcommand{\QuestionTwoSix}{%
Compute the resistive and displacement current densities in the cable.
Is current conservation satisfied? Which current dominates in each material, and why?}
\newcommand{\QuestionTwoSeven}{%
Compute the per-unit-length capacitance (\(\mu\mathrm{F/km}\)) and compare with
the analytical value (13) and with the datasheet. State the assumptions of the
analytical formula (in particular, how the semiconducting layers are treated).
Does the capacitance vary with frequency? Why (not)?}
\newcommand{\QuestionTwoEight}{%
Propose an improvement of the cable design with respect to the electric stress,
and support it quantitatively.}
